% problem-000045 5 铈四价水解与电位
\section*{5. 铈四价水解与电位}
\textit{下列为分问。}

\subsection*{5-1}
尽管⽔溶液中 $\mathrm { C e ^ { 4 + } }$ 强烈⽔解，但在⾼酸度溶液 $( \mathrm { [ H C l O _ { 4 } ] } > 6 . 0 5 \mathrm { m o l } \mathrm { k g ^ { - 1 } } )$ 中 $\mathrm { C e ^ { 4 + } }$ +仍然会发⽣反应： $\mathrm { C e ^ { 4 + } + }$ $\mathrm { e } ^ { - } \stackrel { \right. } { \left. } \mathrm { C e } ^ { 3 + }$ 。⾼氯酸溶液中 $\mathrm { C e ^ { 4 + } }$ 与OH−形成4种配合物的累积⽣成常数的对数分别为lg $\beta _ { 1 } ^ { \ominus } = 1 4 . 7 6 .$ 、lg $\beta _ { 2 } ^ { \ominus } =$ 28.04、lg $\beta _ { 3 } ^ { \ominus } = 4 0 . 5 3$ 、lg $\beta _ { 4 } ^ { \ominus } = 5 1 . 8 6$ ，有关物质的标准摩尔⽣成吉布斯⾃由能 $\Delta _ { \mathrm { f } } G _ { \mathrm { m } } ^ { \ominus }$ 列于表5-1。

表5-1 物质的标准摩尔⽣成吉布斯⾃由能（298.15 K）

\textit{（原卷此处含表格/仪器试剂表，详见 PDF 或 MD 源文件。）}

\subsection*{5-1}
-1 写出 $\mathrm { C e ^ { 4 + } }$ ⽔解反应的⼀般表达式、标准⽔解平衡常数的⼀般表达式。

\subsection*{5-1}
-2 写出 $\mathrm { C e ^ { 4 } }$ +与OH−形成配合物反应的⼀般表达式、标准累计⽣成常数 ${ \mathrm { \it B } } _ { i } ^ { \ominus }$ 的⼀般表达式。

\subsection*{5-1}
-3 求 $\mathrm { C e ^ { 4 + } }$ ⽔解反应⽣成的4种氢氧化物的 $\Delta _ { \mathrm { f } } G _ { \mathrm { m } } ^ { \ominus } .$

\subsection*{5-1}
-4 当 $\mathrm { p H } < 8$ $\mathrm { C e ^ { 3 + } }$ +不⽔解。当 $\mathrm { H C l O _ { 4 } }$ 浓度 $< 6 . 0 5 \mathrm { m o l } \mathrm { k g ^ { - } }$ 1时， $\mathrm { C e ^ { 4 + } / C e ^ { 3 + } }$ 电对的电位是溶液的pH、$\mathrm { C e ( I V ) }$ 配合物的 $\beta _ { i } ^ { \ominus }$ 、有关物种浓度的函数，试导出这⼀函数关系。

\subsection*{5-2}
在 $\mathrm { H C l O _ { 4 } }$ 溶液中， $\mathrm { C e ^ { 4 + } / C e ^ { 3 + } }$ 电对的电位与温度的关系如下图所⽰。

（图：）

根据上图指出，当 $\mathrm { H C l O _ { 4 } }$ 浓度⼤于或⼩于1.85 m时， $\mathrm { C e ^ { 4 + } }$ 被还原成 $\mathrm { C e ^ { 3 + } }$ 的反应的熵变∆�⼤于0、等于0或⼩于0，并说明反应焓变∆�的相应变化情况。

\subsection*{5-3}
采⽤分光光度法研究了⾼氯酸中Ce(IV)对有机物ATF的氧化反应。所有动⼒学测试都在准⼀级条件下进⾏，即ATF的浓度⼤⼤过量于[Ce(IV)]。研究发现，Ce(IV)氧化ATF的反应对Ce(IV)的分级数为⼀级，求得的表观⼀级速率常数 $\smash {  { \mathcal { k } } \not \equiv  { \mathcal { G } } _ { \perp } ^ { } } | |$ 于表5-2。该反应的化学计量反应式如下：

（图：）
表 5-2 [ATF]和[H+]对 Ce(IV)氧化 ATF 的反应的 $| k _ { \mp } \frac { \Xi } { \Xi } \rangle$ 响 (293.3 K， $[ \mathrm { C e } ( \mathrm { I V } ) ] = 2 . 0 { \times } 1 0 ^ { - 4 } \ : \mathrm { m o l } \ : \mathrm { L } ^ { - 1 } )$

\textit{（原卷此处含表格/仪器试剂表，详见 PDF 或 MD 源文件。）}

结果表明，Ce(IV)氧化ATF的反应是⾃由基反应。在题给酸度条件下，Ce(Ⅲ)不⽔解，于是提出如下反应机理：

（图：）

()

()

（图：）

()

（图：）

()

（图：）

()

$
\mathrm{O} \xrightarrow [ \mathrm{NMe} _ {2} ]{\mathrm{OH}} \quad \xrightarrow {\text {fast}} \quad \mathrm{CO} _ {2} + \quad \mathrm{HNMe} _ {2}\tag{6}
$

已知： $K _ { \mathrm { O H } } = 0 . 2 7 \mathrm { m o l } \mathrm { L } ^ { - 1 } , K = 1 2 0 \mathrm { m o l } \mathrm { L } ^ { - 1 } \mathrm { q }$

\subsection*{5-3}
-1 请根据以上反应机理，导出Ce(IV)氧化ATF反应的速率⽅程，并给出 $k \#$ 的表达式。

\subsection*{5-3}
-2 根据 $k \#$ 的表达式，结合表5-2，讨论[ATF]对�_{表}的影响，指出反应对ATF的级数⼤于1或⼩于 $1 _ { \circ }$

\subsection*{5-3}
-3 根 $\ddagger \boxed { \mathrm { E } } k \ne $ 的表达式，结合表5-2，讨论[H+]对 $k _ { ☉ }$ 的影响。
